\section{Introduction}\label{sec:intro}

Does there exist a deterministic multitape Turing machine that transposes every $d \times d$ matrix of $b(d)$-bit entries,
where $b(d) = \max(1, \clog{d})$ and $L = d^2 b(d)$ is the bit-length of the input, in time $O(L)$? Following the earlier
papers of this series~\cite{FirstPaper2026}, we write~(T) for this assertion:
\begin{quote}
(T)\quad There exists a deterministic multitape Turing machine---with a binary read-only two-way input tape, a binary
write-only one-way output tape, and a fixed finite number of two-way read/write work tapes over a fixed finite work
alphabet---that transposes $d \times d$ matrices of $b(d)$-bit entries in $O(d^2 b(d))$ steps.
\end{quote}
By a theorem of Harvey and van der Hoeven~\cite{HvdH2025}, $N$-bit integer multiplication reduces in linear time on a
multitape Turing machine to matrix transposition; hence refuting~(T) would establish an $\Omega(N \log N)$ lower bound
for deterministic off-line multitape integer multiplication. Whether~(T) holds remains open.

\subsection{Context in the series}

The fourth paper of this series~\cite{FourthPaper2026} reduced the impossibility of \emph{oblivious} linear-time
transposers to the conjunction of two combinatorial statements that make no mention of Turing machine states:
\begin{enumerate}[label=(\roman*)]
\item \emph{The trajectory cut problem} ($\TCP(k)$): in every $k$-head walk dependency DAG of length $T$, deleting at
  most $\eps T$ edges reduces the longest directed path to length $O_{\eps,k}(1)$.
\item \emph{The coding crux}: every storage code over a finite alphabet on a transitive-tournament-free (TT-free)
  directed recovery relation of bounded degree has asymptotic rate bounded away from~$1$---and, in its sharpest form
  (Conjecture~P of~\cite{SeventhPaper2026}), rate at most~$1/2$ whenever every vertex has in-degree at most~$3$.
\end{enumerate}
Subsequent papers sharpened both halves of this program while exposing concrete obstructions to naive extensions:
\begin{itemize}
\item On the \emph{directed coding} side, the fourth paper~\cite{FourthPaper2026} proved the half-rate bound for
  \emph{scalar linear} codes on the mixed class $\ClassS$ of TT-free relations with maximum in-degree~$3$ in which every
  vertex of in-degree~$3$ has only reciprocal children (its Theorem~S). The seventh paper~\cite{SeventhPaper2026}
  introduced the nested elimination engine over general polymatroids (removing linearity altogether), proved the
  all-code half-rate bound on all undirected subcubic triangle-free graphs (its Theorem~U3), and constructed the
  $15$-vertex TT-free orientation $\Ex$ of maximum in-degree~$3$ on which minimum-core peeling fails. Within item~2
  of~\cite[Section~9.2]{SeventhPaper2026}, alongside Conjecture~P (\cite[Conjecture~9.1]{SeventhPaper2026}, for general
  in-degree at most~$3$), it singled out as an intermediate target whether the fourth paper's Theorem~S lifts from scalar
  linear codes to \emph{all} codes on the entire class~$\ClassS$, and noted two unformalized impossibility
  claims~\cite[Remarks~7.3 and~7.5; Section~9.2, item~6]{SeventhPaper2026}: that $\Ex$ admits no piece certificate on its
  unique minimum nonempty core over any number of records, and that the Petersen graph admits no global elimination
  certificate over any number of records.
\item On the \emph{degree-$4$ linear coding} side, while at degree~$5$ scalar binary storage codes already exceed
  rate~$1/2$ on the Clebsch graph (rate~$5/8$, Lean-verified in~\cite[Corollary~6.11]{FourthPaper2026}), whether any
  storage code on a triangle-free graph of maximum degree~$4$ can exceed rate~$1/2$ remains open
  (\cite[Section~9.2, item~4]{SeventhPaper2026}; \cite[Section~10, Problem~10.4]{EighthPaper2026}), and tree-local
  arguments cannot handle arbitrary clusters of degree-$4$ vertices~\cite[Proposition~6.28]{FourthPaper2026}. The eighth
  paper~\cite{EighthPaper2026} isolated the exceptional set $H$ of degree-$4$ rows in a unit-diagonal check matrix
  $M = I + B$ on a triangle-free graph of maximum degree~$4$ and proved $\dim\ker M \le \card{V}/2$ over every field
  under two hypotheses on $H$: (S1)~no two vertices of $H$ share a common neighbor, and (S2)~no two vertices of $H$ are
  adjacent. Whether~(S2) is genuinely needed or an artifact of the induction was left
  open~\cite[Section~8.2, Remark~8.8; Section~10, Problem~10.4]{EighthPaper2026}.
\item On the \emph{head-motion and deep-nesting} side, the eighth paper~\cite{EighthPaper2026} proved $\TCP(k)$ for all
  schedules of bounded nesting load (its Theorem~F) and isolated the $2$-head \emph{deep-nesting pair} $\DNP(s_0, m)$,
  running two shifted copies of the recursive nested walk $U(s_0, m, 1)$, as the sharpest open case of $\TCP(2)$. On a
  single copy of $U(s_0, m, 1)$, the eighth paper's turn counter~\cite[Section~5.2, Theorem~5.2 and Corollary~5.5]{EighthPaper2026}
  (built on Lemma~Zq of~\cite[Section~5.5, Lemma~5.15]{FourthPaper2026}) cuts depth to $2^{m + O(1)}/\eps$; the eighth
  paper introduced block-cancellation labels $B(a, K)$ to probe whether two heads can cancel turn-counter drift across
  scales, stated Conjecture~A$''$~\cite[Section~5.6, Conjecture~5.24]{EighthPaper2026}---that at some odd shift, every
  mod-$q$ label of cost at most $n_{\mathrm{h}}/8$ on the $r$-fold cover has at most $C\log_2 q + C' + \eta(m+1)$ doubly
  $\theta_0$-served levels ($\eta < 1$)---and at the end of~\cite[Section~5.6]{EighthPaper2026} identified two open steps
  for proving Conjecture~A$''$ even on block-label pairs: (1)~\emph{``a bound on $\abs{F}$ on the straddling hops, or a
  version of Theorem~V that counts the window directly''}, and (2)~\emph{``at one fixed shift for all coarse classes,
  equidistribution of a class's apex centres relative to the other head's canceller valleys of level about
  $k + \log_2(m - k)$, a mixed-radix carry problem''}. Meanwhile, the ninth
  paper~\cite{NinthPaper2026} developed an independent bypass of $\TCP$ via complete source covers and return-profile
  barriers on restricted return geometries, leaving open the unrestricted adaptive regime.
\end{itemize}

\subsection{Summary of main results}

This paper resolves the intermediate directed coding target of~\cite[Section~9.2, item~2]{SeventhPaper2026} (while
Conjecture~P, its Conjecture~9.1, remains open), removes hypothesis~(S2) from Theorem~SH of~\cite{EighthPaper2026}, and
proves and formalizes structural stroke, straddle, and super-Fibonacci laws together with three-regime failure witnesses for
bounded-range block-label pairs on the deep-nesting walk. Every new numbered result in the paper is verified in Lean~4 using
only the three standard foundational axioms (\cref{sec:prelim-lean,app:verify}).

\paragraph{1.\ Directed reciprocity, digon-cycle dichotomy, and the all-code Theorem~S$^*$ (\cref{sec:reciprocity}).}
We begin by recasting the stateful recursive predicate \lean{HalfRate.RelCharged} of~\cite{SeventhPaper2026} in purely
positional terms (\cref{prop:PN}): an elimination list $l = [v_0, \dots, v_{k-1}]$ with records $\{P_r\}_{r \in I}$,
supplying-record map $b$, and optional outside-parent map $\mathrm{par}$ is relationally nested from empty initial state
if and only if every step covers its active parents from the revealed prefix $S_j$, its supplying record $P_{b(v_j)}$,
and its optional outside parent, and whenever $j < j'$ share the same outside parent $w$, $P_{b(v_j)} \subseteq S_{j'}
\cup P_{b(v_{j'})}$ (for disjoint records, either $b(v_j) = b(v_{j'})$ or all of $P_{b(v_j)}$ is revealed before step
$j'$). As an immediate invariant on any piece certificate, no two vertices joined by a reciprocal pair (a \emph{digon})
can belong to the same own record $P_{a(u)}$, so every connected digon component on at least two vertices requires
$m \ge 2$ records.

We then classify piece certificates on chordless digon cycles $C_L$ ($L \ge 4$) in TT-free relations where each cycle
vertex has at most one active parent outside $U \cup C_L$ (\cref{prop:CS,thm:DC}): an $m$-record piece certificate
exists on $C_L$ for \emph{some} $m \ge 1$ if and only if the outside parents do not form a \emph{stride-$3$ hub triple},
that is, $3 \mid L$ together with outside vertices $w_0, w_1, w_2 \notin U \cup C_L$ such that $w_{i \bmod 3}$ is an
active parent of cycle vertex $i$ for every $0 \le i < L$. When the stride-$3$ hub obstruction is absent, $m = 2$
records suffice for $L = 4$ and $m = L + 1$ records suffice via two-sided arc expansion for every $L \ge 5$ (while on
pure stride-$s$ hub cores $C_{sk}$, $m = 2$ suffices for even $s \ge 4$ and $m = 3$ suffices for every $s \ge 4$); when
the stride-$3$ hub triple is present, a positional locking invariant rules out piece certificates across \emph{all}
$m \ge 1$. Specializing to the unique minimum nonempty core $\{0,\dots,5\}$ of $\Ex$ (\cref{cor:x15}) and to generalized
Petersen graphs (\cref{prop:GP}) resolves both parts of~\cite[Remarks~7.3 and~7.5; Section~9.2, Problem~6]{SeventhPaper2026}
in Lean~4 and reveals
a non-monotone global-certificate record hierarchy: $\GP(5,2)$ (the Petersen graph) requires $\infty$ records globally,
$\GP(7,2)$ requires exactly $4$ records, and $\GP(9,2)$ drops back to $3$ records.

Using these cycle tools, we prove \textbf{Theorem~S$^*$} (\cref{thm:Sstar}): every relation $N$ in the full class
$\ClassS$ of Theorem~S of~\cite{FourthPaper2026} admits a certified piece decomposition (\lean{HalfRate.Decomp}), so
every valid polymatroid $h$ on $V$ obeys $2 h(V) \le \sum_{v \in V} h(\{v\})$ and every
storage code $C \subseteq \alpha^V$ (linear or nonlinear, over any finite alphabet $\alpha$) satisfies
$\card{C}^2 \le \card{\alpha}^{\card{V}}$. The key mechanism overcoming embedded stride-$3$ hub cycles inside an
in-closed strongly connected component $S$ of $\ClassS$ is \emph{hub augmentation}
(\lean{exists_piece_of_scc_classS}): if $S$ has no minimal chordal digon arc and a shortest digon cycle $C_L \subseteq S$
is obstructed as a standalone piece, then $C_L$ is a stride-$3$ hub cycle with hubs $w_0, w_1, w_2 \in S \setminus C_L$;
strong connectivity forces an outgoing arc from some $u_{i_0} \in C_L$ to $S \setminus C_L$, which $\ClassS$ reciprocates
into a digon between $u_{i_0}$ and its hub $w_{r_0}$ while girth minimality forces $w_{r_0}$ to have in-degree~$2$ and no
other child on $C_L$, so adjoining $w_{r_0}$ yields an $(L+2)$-record sequential nested piece certificate on
$C_L \cup \{w_{r_0}\}$. We deduce a directed vertex-cut and isolated-outside-vertex barrier inequality for all $\ClassS$
relations (\cref{cor:Bstar}) and construct an explicit $21$-vertex strongly connected triangle-free witness $\Gto \in \ClassS$
(\cref{prop:G21}) on which low-in-degree peeling, one-way cycle peeling, and shortest-digon-cycle rules all fail
simultaneously, while Theorem~S$^*$ certifies rate at most~$1/2$.

\paragraph{2.\ Scalar linear codes at degree~$4$ without heavy separation (\cref{sec:crux}).}
In \cref{sec:crux} we prove \textbf{Theorem~SH$^*$} (\cref{thm:SHstar}): over any field $F$, if $M \in F^{V \times V}$
has unit diagonal, no recovery triangle, and satisfies $(\star 4)$ at every vertex (for instance when $M$ is supported in a
triangle-free simple graph $G = (V, E)$ of maximum degree~$4$), and the heavy set $H_M$ of vertices violating $(\star\star)$
either has size at most~$2$ or satisfies~(S1) alone---no vertex has two distinct heavy recovery neighbors (in graph form,
no two distinct vertices of $H_M$ share a common neighbor in~$G$)---then $\dim_F \ker M \le \lfloor \card{V}/2 \rfloor$
(and likewise under the same hypotheses on $M^\T$). Hypothesis~(S2) of~\cite{EighthPaper2026} is eliminated
completely: $G[H_M]$ may contain arbitrary induced matchings of adjacent degree-$4$ vertices. The proof replaces Phase~E
of~\cite{EighthPaper2026} with a unified final count (\lean{F3_star}): under~(S1) alone, a pending triple of a heavy vertex
$h$ contains at most one heavy member $g$, leaving two light members $u_1, u_2$ whose shared pending pair $\{a, b\}$ and
neutral dimension drop contradict linear independence uniformly for all $\card{W \setminus X_0} \ge 3$, while in minimal
counterexamples any degree-$4$ heavy vertex still retains at least three light neighbors under~(S1). Combining
Theorem~SH$^*$ with Tutte--Berge deficiency yields \textbf{Corollary~M$^*$} (\cref{cor:Mstar}), improving the eighth
paper's distance-$3$ packing invariant $\rho_3(G, H_M)$ to the common-neighbor-free packing invariant
$\rhom(G, H_M) \ge \rho_3(G, H_M)$.

\paragraph{3.\ Deep-nesting stroke laws, super-Fibonacci unwrapping, and bounded-range block pairs (\cref{sec:dnp}).}
In \cref{sec:dnp} we formalize the algebraic and combinatorial theory of turn-supported periodic labels modulo~$q$ on the
deep-nesting walk $U(s_0, m, 1)$ and on the deep-nesting pair $\DNP(s_0, m)$ (thereby also providing Lean~4
counterparts for the unformalized stroke, band, block and straddle lemmas of~\cite[Section~5]{EighthPaper2026}):
\begin{itemize}
\item \textbf{Proposition~SL} (\cref{prop:SL}) establishes the exact stroke decomposition bijection in $[0, \Pi_0)$,
  proves that under fair binary branch choices the stopping depth before the nested walk visits~$m$ has distribution
  $2^{-(s+1)}$ on $\{0,\dots,m-1\}$ independently of directional parity, and constructs an $O(m)$ recursive turn-potential
  evaluator \lean{evalPhiRec} proved equal to cumulative turn jumps over a full period $P = 4 \Pi_0$.
\item \textbf{Theorem~V$^*$} (\cref{thm:Vstar}) proves an $m$-free mod-$q$ straddle involution bound: whenever a level-$k$
  apex hop of one head straddles a level-$\kappa$ V-turn of the other head ($\Pi_\kappa \le R_k$), the paired half-widths
  $(d, \Pi_\kappa - d)$ cancel the nested potential identically over $\Z$, forcing their reduced mod-$q$ errors to sum to
  at least $(2 A' + 2(c_{\sigf}(\kappa) - c_{\sigb}(\kappa))) \bmod q$; moreover, summing the straddle set size over
  centers of either fixed parity across a $2\Pi_\kappa$-window yields the \emph{exact} triangular identity
  $(\Pi_\kappa / 2)^2 = \Pi_\kappa^2 / 4$.
\item \textbf{Theorem~LU} (\cref{thm:LU}) proves that for any level-uniform turn-supported label modulo $q \ge 3$, if
  $s \ge 1$ consecutive levels have single-head reduced apex and valley band-$0,1,2$ errors at most $q/8$, the modular
  band relation unwraps over $\Z$ with zero wrap-around and forces the super-Fibonacci inequality
  $b_0(\ell) \ge b_0(\ell+1) + b_0(\ell+2) + 4$, whence $32(F_s - 1) \le q$ and $s \le 2\lfloor \log_2 q \rfloor + 2$.
\item \textbf{Proposition~OC} (\cref{prop:OC}) proves exact forward and backward block-cancellation bounds for $B(a, K)$,
  establishes strictly positive mean drift $+(a(K+1)+1)/2$ per $(a+1)$-level block across all $2^{a+1}$ stroke-direction
  choices, and shows that the pure binary clock $B(\infty, 0)$ fails modulo every Mersenne number $q = 2^B - 1$ ($B \ge 2$)
  with straddle residue $q - 1$ at every third/periodic scale $B \mid (m - \kappa + 2)$.
\item \textbf{Theorem~A} (\cref{thm:A}) combines these ingredients for block-label pairs $(B(a_1, K_1),\allowbreak B(a_2, K_2))$ of
  bounded block length ($1 \le a_1, a_2$): Theorem~V$^*$ completely bypasses open step~(1) of~\cite[Section~5.6]{EighthPaper2026}
  for level-uniform labels (paired straddle sums need no bound on $\abs{F}$ and no window assumption on~$q$), while
  \lean{exists_oddShift_coarse_straddle} establishes a one-direction $\tfrac{1}{16}$-fraction form of step~(2) (weaker than
  the all-coarse-classes hypothesis of~\cite[Proposition~5.23]{EighthPaper2026}, which is obstructed by prefix lockstep of
  nested V-apices, \lean{apexOffset_diff_indep_prefix}). Whenever $m < s_0$ and
  $4(\lfloor \log_2(m+1)\rfloor + 1 + \max(a_1, a_2)) \le m + 1$, Theorem~A produces two fixed odd shifts $\xi_1, \xi_2$
  (a single odd shift $\xi$ for symmetric pairs $a_1 = a_2, K_1 = K_2$) giving three-regime failure witnesses on
  $\Omega(m)$ levels across \emph{every} $K_1, K_2 \in \Z$ and \emph{every} modulus $q \ge 3$, with \emph{zero} $\log q$
  loss: blind single-head error witnesses in Regimes~1 and~2, and two-head coarse straddle witnesses in Regime~3. We also
  record the four remaining steps separating these witnesses from Conjecture~A$''$ (\cref{rem:scope-A}).
\end{itemize}

\subsection{Related work}

The equivalence between locally recoverable storage capacity on a digraph $G$, index coding broadcast rate on the
complement, and the guessing number $\gn(G) = \sup_{\card{\alpha} \ge 2} \sup_{C} \log_{\card{\alpha}} \card{C}$ (over
locally recoverable codes $C \subseteq \alpha^V$,~\cite{Riis2007,SeventhPaper2026}) goes back to Birk and Kol~\cite{BirkKol2006},
Riis~\cite{Riis2007}, Bar-Yossef, Birk, Jayram and Kol~\cite{BBJK2011}, Gadouleau and Riis~\cite{GadouleauRiis2011}, and
Mazumdar~\cite{Mazumdar2014}. On undirected triangle-free graphs, Christofides and Markstr\"om~\cite{ChristofidesMarkstrom2011}
bounded the guessing number by the fractional matching number (giving rate~$1/2$ for all undirected triangle-free graphs
only on symmetric recovery neighborhoods), Cameron, Dang and Riis~\cite{CameronDangRiis2016} classified high-guessing
triangle-free graphs in low degree, Atkins, Rombach and Skerman~\cite{AtkinsRombachSkerman2017} determined guessing
numbers of odd cycles, and Mazumdar, McGregor and Vorotnikova~\cite{MMV2019} proved half-rate bounds under sparse global
average degree, whereas without a degree bound Barg and Z\'emor~\cite{BargZemor2022} and Barg, Schwartz and
Yohananov~\cite{BSY2024} constructed coset and Cayley storage codes on triangle-free graphs of growing degree with rates
exceeding~$1/2$ (approaching~$3/4$ and~$1$, respectively).
On directed graphs without undirected symmetry, undirected fractional-matching duality breaks down because of directed
cycles of pure one-way arcs; Theorem~S$^*$ (\cref{thm:Sstar}) is, to our knowledge, the first all-code half-rate theorem
covering a class of TT-free digraphs that simultaneously permits full subcubic undirected reciprocity, unbounded
directed one-way cycles, and mixed components combining both.

On the complexity side, reductions from network coding and circuit depth to Valiant's edge-removal problem go back to
Valiant~\cite{Valiant1977}; see~\cite{FirstPaper2026,FourthPaper2026,EighthPaper2026,NinthPaper2026} for a detailed
account of the connections between $\TCP(k)$, monotone and non-monotone depth reduction, matrix rigidity, and source
covers.

\subsection{Organization and formal reading guide}

\Cref{sec:prelim} fixes notation from the earlier papers of the series and explains our Lean~4 verification conventions
(\cref{sec:prelim-lean}). \Cref{sec:reciprocity} develops positional nestedness, the digon-cycle dichotomy, the
generalized Petersen record hierarchy, and Theorem~S$^*$. \Cref{sec:crux} proves Theorem~SH$^*$ and Corollary~M$^*$.
\Cref{sec:dnp} proves Proposition~SL, Theorem~V$^*$, Theorem~LU, Proposition~OC, and Theorem~A on the deep-nesting walk
and pair. \Cref{sec:construction} synthesizes what any hypothetical linear-time transposer must look like in light of
these bounds, and \cref{sec:open} collects the sharpened open problems. \Cref{app:index} indexes every Lean declaration
cited in the paper, and \Cref{app:verify} gives exact instructions for building the library, checking axioms, and
reproducing the finite numerical examples.
